%
% connection_lease_reaping.tex
%
% A Z specification of the Hub's client-connection lease/reap lifecycle and
% its coupling to scene and element ownership.
%
% REFINEMENT (round 2, bead lux-d84d). Round 1 of this model collapsed every
% departure into one atomic Reap operation, guarded on the lease having
% lapsed behind a dead transport. That abstraction type-checked and
% model-checked clean, and the implementation built from it passed the full
% suite -- but a four-way review of the shipped code found three ownership-
% integrity holes the single-Reap abstraction could not see, because the
% code does not have one departure path: it has three, and only one of them
% was ever wired to release ownership. This round replaces the single Reap
% with the code's actual shape -- a non-destructive read, a departure
% coordinator, and a write path that must renew its own caller before it
% sweeps anyone else -- and adds the two invariants (I5, I6) that state, in
% general form, why each hole is now closed.
%
% ADDENDUM within round 2: the first cut of this round made every departure
% pull-based -- discharged only by another connection's Claim or an external
% Depart signal. The leader ruled that pull-based reaping cannot satisfy the
% bead's own premise ("the 30s lease exists for this and is NOT firing" --
% the connection persisted 25.79h with nobody else ever writing or
% disconnecting it). TimedReap is added below: the same atomic
% deregister-and-release step, but guarded only on a lapsed lease, so the
% lease is a genuine background backstop and I1 no longer depends on any
% other party's activity.
%
% ROUND 3 (bead lux-vvmt). Design: docs/architecture/lease-reap-cascade.md,
% ratified by the operator. Round 2 modeled one leg of departure --
% registry membership and scene ownership. The shipped code's departure
% cascade is wider: a graceful Depart, via lifecycle.py, also drops the
% connection's Hub subscriptions and writer binding (Hub.on_disconnect)
% and fires its transport-owned sink (inbox.drop_session); TimedReap and
% Claim's embedded sweep never reached either leg, so a timer-reaped or
% swept-aside connection kept its subscriptions, writer binding, and inbox
% queue forever. Worse, because connection_for derives a STABLE id from a
% client's declared identity, a same-identity reconnect often reuses the
% identical CONN value, so an unreleased writer binding silently blocks a
% legitimate reconnect's own ensure_writer call from installing a fresh
% one -- not a leak, a reconnecting client inheriting a stranger's state.
% This round adds three new state components (hasSubs, hasWriter,
% inboxNonEmpty) and I7, the general form of "leaving the registry
% releases the FULL cascade, however triggered." Grounding the fix
% surfaced a second hazard: register_client/identify_client (Connect/Renew
% below) do not take the Hub's StoreLock, so a same-identity reconnect can
% land inside the window between TimedReap's registry removal and its
% (new) cascade tail, and have its fresh writer binding retroactively
% wiped out by the stale reap's tail. Closing that hazard needs an
% explicit mutual-exclusion device this round adds: TimedReap is split
% into TimedReapBegin/TimedReapEnd around a new cascadeFor flag, and
% Connect/Renew gain a precondition against it -- the formal counterpart
% of bringing register_client/identify_client under StoreLock.write().
% Depart and Claim need no such split: neither is implicated by any
% reported race, and their existing round-2 atomicity already forecloses
% one (see the Invariants section's discussion of why the split is scoped
% to TimedReap alone). I7's statement is scoped accordingly: it excludes
% a connection currently inside its own TimedReapBegin/End window, which
% is not a weakening but the model's honest acknowledgement that nothing
% can observe a real critical section from the inside -- the same
% reasoning round 2 already gives informally for why no operation here
% needs to represent being "mid-Depart" or "mid-Claim."
%
% Read first (round 3, in addition to the round-2 list below): the design
% document's own "Read first" grounding at
% docs/architecture/lease-reap-cascade.md Section 1 and Section 4 --
% domain/hub/hub_display.py (register_client, identify_client,
% drop_connection, _depart_lapsed, apply), domain/hub/lifecycle.py
% (disconnect_connection), domain/hub/hub.py (on_disconnect), and
% domain/hub/inbox.py (ensure_writer, drop_session).
%
% Read first: src/punt_lux/domain/hub/hub_clients.py (HubClientRegistry --
% named_sessions/live_sessions/repos as READS that must become
% non-destructive; reap/_sweep_locked; record/_renewed; discard),
% src/punt_lux/domain/hub/hub_display.py (apply, drop_connection,
% _reap_and_release), src/punt_lux/domain/hub/owner_tracker.py, and the
% graceful-departure path src/punt_lux/session_cleanup.py ->
% src/punt_lux/domain/hub/lifecycle.py's disconnect_connection ->
% HubDisplay.drop_connection, plus src/punt_lux/domain/hub/session_lease.py.
%
% Type-check:  fuzz docs/connection_lease_reaping.tex
% Model-check (see the Verification section for the exact goal predicates):
%   probcli docs/connection_lease_reaping.tex -model_check -p DEFAULT_SETSIZE 3
%
% Formatting follows the fuzz house rules: schema lines under 80 columns,
% predicates broken at \land/\lor/\implies.
%

\documentclass[a4paper,10pt,fleqn]{article}
\usepackage[margin=1in]{geometry}
\usepackage{fuzz}
\usepackage[colorlinks=true,linkcolor=blue,citecolor=blue,urlcolor=blue]{hyperref}

\begin{document}

\title{lux Connection-Lease Reaping: a Z Specification}
\author{Formal model of the Hub client registry's lease sweep and its
  coupling to scene ownership -- refined to the code's multiple departure
  paths}
\date{September 2026}
\maketitle

% ============================================================
\section{Introduction}
% ============================================================

Round 1 of this model proved a single property: once a connection's lease
lapses, one atomic \texttt{Reap} both drops it from the registry and
releases every scene it owned. That property type-checked, model-checked
clean against all four of its invariants, and the implementation built
from it passed 4{,}730 tests and an independent re-verification. It was
also too narrow to be true of the code it claimed to model. The shipped
system does not have one departure operation; it has three, and the
model's single \texttt{Reap} silently assumed all three were the same
thing.

A four-way review of the implementation found three separate places where
a connection can leave the registry without its ownership being released
-- three holes the single-\texttt{Reap} abstraction hid because it never
represented more than one way to leave:

\begin{description}
  \item[H1 --- destructive-read-sweep.] \texttt{HubClientRegistry
    .named\_sessions} (and, through it, \texttt{live\_sessions} and
    \texttt{repos}) calls \texttt{\_sweep\_locked} on every read, which
    drops a lapsed session from the registry's own dict and returns what
    it dropped -- but six call sites read through \texttt{named\_sessions}
    /\texttt{live\_sessions} without ever forwarding that return value to
    \texttt{OwnerTracker}. A read that happens to run after a connection's
    lease has lapsed silently deregisters it, ownership untouched. The one
    call site that \emph{does} forward the return value
    (\texttt{HubDisplay.\_reap\_and\_release}, via \texttt{HubClientRegistry
    .reap}) is a write path, not a read.
  \item[H2 --- graceful-departure-without-release.] \texttt{HubDisplay
    .drop\_connection} -- the ordinary path an MCP session's clean
    disconnect or the SDK's own idle-timeout takes, via
    \texttt{session\_cleanup.py}'s teardown leg and
    \texttt{lifecycle.disconnect\_connection} -- calls
    \texttt{HubClientRegistry.discard} unconditionally and never touches
    \texttt{OwnerTracker} at all. Round 1's fidelity control
    (\texttt{connection\_lease\_reaping\_no\_release\_buggy.tex}) already
    modeled ``no release on departure,'' but it modeled that as a variant
    of the \emph{lease-lapse} path. It is not: it is this path, the one
    every ordinary session end takes, lapsed or not.
  \item[H3 --- self-reap-without-renew.] \texttt{HubDisplay.apply} calls
    \texttt{self.\_reap\_and\_release()} -- an unconditional sweep of
    \emph{every} lapsed session, itself included -- before it does
    anything else, and never calls anything that renews the calling
    connection's own lease. A connection that only ever contacts the Hub
    to patch an already-shown scene never renews outside that first
    \texttt{show}; once its lease's length has passed, its own next write
    sweeps and releases itself before performing the very write that
    should have kept it alive.
\end{description}

All three share one root: the code has no single place where ``leave the
registry'' and ``release what you owned'' are the same atomic step for
\emph{every} trigger. The operator's ruling fixes that with three
decisions this model encodes and does not re-open:

\begin{enumerate}
  \item \textbf{One coordinator.} The only way a connection leaves the
    registry is one atomic deregister-and-release step. Reads never
    remove anyone. Every departure trigger -- graceful disconnect, the
    SDK's idle-reap, and a lease lapsing -- funnels through that one step.
  \item \textbf{Release ownership, keep content.} Departure releases a
    connection's scenes to \texttt{unowned}; it does not tear the scenes
    down. A reconnecting client reclaims through the ordinary
    unowned-claim path -- I4, unchanged.
  \item \textbf{Renew on contact, including a write.} Any authenticated
    contact renews the caller's own lease, and that now explicitly
    includes the ownership-gated write path (\texttt{apply}), which round
    1 left uncovered.
\end{enumerate}

The first cut of this round realized ``a lease lapsing'' entirely as a
\emph{side effect} of some other party's activity -- \texttt{Depart}'s
external signal, or \texttt{Claim}'s embedded sweep of others. The leader's
ruling on that cut's own surfaced gap adds a fourth decision:

\begin{enumerate}
  \setcounter{enumi}{3}
  \item \textbf{The lease is a real backstop, not just a bookkeeping
    field.} A lapsed lease must eventually discharge on its own, with no
    read, no other connection's write, and no detected disconnect required
    -- exactly what ``the 30-second lease exists for this'' promises, and
    exactly what a purely pull-based design cannot deliver for a killed
    process on an otherwise idle Hub.
\end{enumerate}

Eight operations realize this: \texttt{Connect} and \texttt{Renew} are
unchanged from round 1. \texttt{TransportDies} is unchanged.
\texttt{LeaseLapse} drops the transport-death precondition it carried in
round 1 -- H3 is exactly the case of a lease lapsing under a transport
that never dies, and round 1's guard could not represent that case at
all. \texttt{Read} is new: it is the non-destructive query round 1 never
modeled, because round 1 never needed to distinguish a read from anything
else. \texttt{Depart} is the single coordinator -- one schema, both for a
graceful disconnect and for the SDK's idle-reap, since both concretely
call the same \texttt{drop\_connection} today and neither depends on
lease state. \texttt{Claim} models \texttt{apply}: it renews the caller's
own lease first, then sweeps every \emph{other} lapsed connection and
releases what each owned, then performs the ownership-gated write against
the result -- one atomic step, matching \texttt{apply}'s own docstring,
``reaps and releases first, before the ownership check.'' \texttt{TimedReap}
is this addendum's addition: the identical atomic deregister-and-release
predicate \texttt{Depart} and \texttt{Claim}'s sweep already perform,
guarded only on \texttt{lease~c? = llapsed} -- no other connection, no read,
and no external signal is a precondition of it firing.

Four properties survive from round 1, one strengthened:

\begin{description}
  \item[I1 --- transport-gone implies \emph{unconditionally} eventually
    reaped.] Strengthened by this addendum. A connection whose transport
    has died cannot renew, and nothing can indefinitely block
    \texttt{TimedReap} once its lease lapses -- not even the presence of
    every \emph{other} connection idle, every read absent, and no
    disconnect ever detected. Round 1's version of this property read
    ``nothing can indefinitely block \texttt{Depart} once it applies,''
    which is true but was silently relying on \emph{something} eventually
    triggering \texttt{Depart} or a peer's \texttt{Claim} -- exactly the
    dependency the bead's own evidence (a connection persisting 25.79
    hours with nobody else ever writing or disconnecting it) shows cannot
    be assumed.
  \item[I2 --- reap releases ownership.] Once a connection has left
    \texttt{registered} it owns nothing.
  \item[I3 --- at most one live owner per identity.] A live connection is
    never blocked from claiming a scene by a dead predecessor of the same
    declared identity that still, on paper, owns it.
  \item[I4 --- reconnect inherits nothing special (settled).] Unchanged
    from round 1; \texttt{Connect}'s definition has not changed.
\end{description}

Two are new as of the multi-path refinement, and remain the reason that
round exists:

\begin{description}
  \item[I5 --- leaving the registry iff released, however triggered.] The
    general form of I2: no reachable state has a connection absent from
    \texttt{registered} while it still owns a scene, \emph{regardless of
    which operation removed it}, and no read operation ever removes
    anyone. This is what closes H1 and H2 -- H1 broke it via a read
    removing without releasing, H2 broke it via a graceful departure
    removing without releasing; I5 rules out both shapes at once, not
    just the lease-lapse shape I2 alone would cover.
  \item[I6 --- a live writer keeps its ownership.] A connection that keeps
    making authenticated contact by writing is never reaped, by its own
    write or another's, while its transport stays up. This is what closes
    H3: it is I5's negation specialized to the case where no transport
    ever died, so the only way to reach it would be a write path that
    swept its own caller.
\end{description}

This round adds a seventh:

\begin{description}
  \item[I7 --- departure clears the full cascade, however triggered.] The
    general form of I5, extended to the three cascade legs round 2 never
    modeled: no connection outside \texttt{registered}, and not currently
    inside its own \texttt{TimedReapBegin}/\texttt{TimedReapEnd} window,
    still shows a live subscription, a bound writer, or a non-empty
    inbox. This is what closes the gap Section~1 of the design document
    names: \texttt{Depart} already reached all four legs of the cascade
    via \texttt{lifecycle.py}; \texttt{TimedReap} and \texttt{Claim}'s
    embedded sweep reached only the first two (registry, ownership),
    leaving a timer-reaped or swept-aside connection's subscriptions,
    writer binding, and inbox queue live forever.
\end{description}

Ten operations realize round 3 (up from round 2's eight).
\texttt{Connect}, \texttt{TransportDies}, \texttt{LeaseLapse}, \texttt{Read},
\texttt{Depart}, and \texttt{Claim} keep round 2's shape, each carrying the
three new state components through unchanged (or, for \texttt{Depart} and
\texttt{Claim}, gaining three more \(\oplus\) terms in exactly the predicate
that already zeroes \texttt{owner} -- textually the same kind of change both
operations already carry from round 1 to round 2). \texttt{Renew} keeps its
effect but gains a new guard. \texttt{Activate} is new: it models, as one
abstract step, an already-connected client subscribing to a topic, having a
writer bound, or writing to its own inbox -- the design document's own
framing, "modeled abstractly," since this round's invariant needs only that
these three facts can become true, not which topic or which message.
\texttt{TimedReap} is retired as a single schema and replaced by
\texttt{TimedReapBegin}/\texttt{TimedReapEnd}, a two-step decomposition of
what was previously one atomic predicate. This is the one place this round
deviates from the design document's literal sketch, and the deviation is
argued in full where \texttt{TimedReapBegin} is introduced below: reproducing
the register\_client-vs-reap race the design document reports (its
Section~4) is a mathematical impossibility under an unsplit, fully atomic
\texttt{TimedReap}, because a Z operation's guard is evaluated fresh against
the current state at the moment it is chosen to fire, and \texttt{Connect}
always renews the lease it touches -- so no ordering of two atomic,
guard-checked operations can ever place a stale reap decision's effect after
a fresh reconnect's effect on the same connection. The design document itself
says as much without naming the consequence: "the two operations' effects
apply out of the order their real-world causes occurred in" is a
non-sequentially-consistent execution, and a non-sequentially-consistent
execution cannot be witnessed by any total ordering of whole, atomic steps.
Splitting \texttt{TimedReap} at the exact point its real-code counterpart
(\texttt{\_depart\_lapsed}) already splits -- registry-and-ownership removal
first, cascade tail second, with nothing atomic connecting the two once
\texttt{register\_client} stops sharing the lock -- is what makes the
hazard, and the fix, formally checkable at all.

% ============================================================
\section{Basic Types}
% ============================================================

\begin{zed}
  [CONN, IDENT, SCENE]
\end{zed}

Unchanged from round 1. Three \texttt{CONN} values suffice to model an old
and a new connection sharing one identity, plus one connection of a
second identity, plus (for the H3 trace) a distinct claimant. Two
\texttt{IDENT} values exhibit the same-identity and cross-identity cases
side by side. Two \texttt{SCENE} values show one scene surviving a
departure intact beside one that does not. All three carriers are bounded
by ProB's \texttt{DEFAULT\_SETSIZE}.

% ============================================================
\section{Free Types}
% ============================================================

\begin{zed}
  TSTATE ::= tup | tdown
\end{zed}

Whether a connection's transport is still alive, unchanged from round 1.

\begin{zed}
  LSTATE ::= llive | llapsed
\end{zed}

A connection's lease, abstracted to whether it has lapsed, unchanged from
round 1 in its two values -- but see \texttt{LeaseLapse} below for what
changed in when it may transition.

\begin{zed}
  OWNER ::= unowned | conn \ldata CONN \rdata
\end{zed}

Unchanged from round 1.

\begin{zed}
  FSTATE ::= fclear | fset
\end{zed}

New in round 3. A coarse two-value carrier for the three new cascade-leg
flags below, in the same style as \texttt{TSTATE} and \texttt{LSTATE} --
prefixed constructors, not \fuzz's logical \texttt{true}/\texttt{false}, so a
reader is never tempted to read \texttt{hasSubs~c? = fset} as a bare
proposition.

% ============================================================
\section{State}
% ============================================================

Unchanged from round 1: one flat schema, structural coherence only, the
safety properties deliberately absent so a violating operation is a
counterexample rather than a disabled transition.

\begin{schema}{ConnReg}
  registered : \power CONN \\
  identity : CONN \pfun IDENT \\
  transport : CONN \pfun TSTATE \\
  lease : CONN \pfun LSTATE \\
  owner : SCENE \fun OWNER \\
  hasSubs : CONN \fun FSTATE \\
  hasWriter : CONN \fun FSTATE \\
  inboxNonEmpty : CONN \fun FSTATE \\
  cascadeFor : \power CONN
\where
  registered \subseteq \dom identity \\
  \dom identity = \dom transport \\
  \dom transport = \dom lease \\
  \# cascadeFor \leq 1
\end{schema}

As in round 1, \texttt{identity} is kept for every connection that has
ever connected, including a departed one -- I3's same-identity comparison
needs it, and no operation below ever shrinks \texttt{dom identity}.
\texttt{owner} stays total over \texttt{SCENE} for the same reason round 1
gave: every predicate below is then a plain function application, never a
domain-membership guard. The three new cascade-leg flags are total over
\texttt{CONN} for the identical reason -- round 1's argument for
\texttt{owner} applies verbatim, and it sidesteps a real modeling trap: a
partial function keyed like \texttt{lease} would need every operation that
introduces a fresh connection id to decide, separately, whether that id is
new (needing a fresh \texttt{fclear} entry) or a reconnect (needing its
prior value preserved) -- an \emph{if-then-else} on domain membership this
model has no need to carry. \texttt{cascadeFor} is new: a set of at most
one connection, standing for whichever connection, if any, is currently
between \texttt{TimedReapBegin} and \texttt{TimedReapEnd} -- the model's
one deliberately-exposed window onto what is, in the fixed real system, a
single unbroken hold of \texttt{StoreLock}.

% ============================================================
\section{Initialization}
% ============================================================

\begin{schema}{Init}
  ConnReg'
\where
  registered' = \emptyset \\
  identity' = \emptyset \\
  transport' = \emptyset \\
  lease' = \emptyset \\
  owner' = (\lambda s : SCENE @ unowned) \\
  hasSubs' = (\lambda c : CONN @ fclear) \\
  hasWriter' = (\lambda c : CONN @ fclear) \\
  inboxNonEmpty' = (\lambda c : CONN @ fclear) \\
  cascadeFor' = \emptyset
\end{schema}

The first five lines are unchanged from round 1. The four new lines start
every cascade-leg flag \texttt{fclear} and no cascade in flight -- the only
sensible empty state, since nothing has connected yet.

% ============================================================
\section{Operations}
% ============================================================

\subsection{Connect --- arrival and reconnect, undifferentiated (I4)}

Unchanged from round 1: no precondition, so nothing here privileges a
fresh connection over a reconnecting one, and nothing here evicts or
transfers another connection's state. This absence of a guard remains the
whole of I4's argument.

\begin{schema}{Connect}
  \Delta ConnReg \\
  c? : CONN \\
  i? : IDENT
\where
  cascadeFor = \emptyset \\
  registered' = registered \cup \{ c? \} \\
  identity' = identity \oplus \{ c? \mapsto i? \} \\
  transport' = transport \oplus \{ c? \mapsto tup \} \\
  lease' = lease \oplus \{ c? \mapsto llive \} \\
  hasWriter' = hasWriter \oplus \{ c? \mapsto fset \} \\
  hasSubs' = hasSubs \\
  inboxNonEmpty' = inboxNonEmpty \\
  cascadeFor' = cascadeFor \\
  owner' = owner
\end{schema}

Two changes from round 1/2. First, the guard: \texttt{$cascadeFor = \emptyset$}
is new -- the formal counterpart of bringing \texttt{register\_client} under
\texttt{StoreLock.write()} (design document Section~4). Second, the effect:
\texttt{$hasWriter' = hasWriter \oplus \{c? \mapsto fset\}$} models
\texttt{inbox.ensure\_writer}'s net effect on any (re)connect -- either a
genuinely fresh writer is bound because none existed, or one already did and
\texttt{ensure\_writer} returns early, leaving it \texttt{true} either way;
this coarse model has no way to distinguish a fresh writer from a stale one
by identity, so it does not need to. \texttt{hasSubs} and
\texttt{inboxNonEmpty} are left unchanged by \texttt{Connect} -- deliberately:
a reconnect does not automatically resubscribe to any topic, and it does not
automatically drain an inbox queue that already had messages in it. This
asymmetry is not an oversight; it is exactly Section~1.1 of the design
document's own diagnosis, that a reconnecting client silently inherits
whatever \texttt{hasSubs}/\texttt{inboxNonEmpty} state a departed
predecessor under the same identity left behind, unless departure itself
cleared it first.

\subsection{Renew --- authenticated contact that is not a write}

Models \texttt{HubClientRegistry.record} called without an
\texttt{AddElement}/\texttt{SetProperty}/\texttt{RemoveElement} attached
to it -- an explicit \texttt{identify} or \texttt{register\_client} call.
Guarded on \texttt{transport~c? = tup}, unchanged from round 1: a
connection with a dead transport cannot make contact of any kind, so it
cannot renew this way either.

\begin{schema}{Renew}
  \Delta ConnReg \\
  c? : CONN
\where
  c? \in registered \\
  transport~c? = tup \\
  cascadeFor = \emptyset \\
  lease' = lease \oplus \{ c? \mapsto llive \} \\
  registered' = registered \\
  identity' = identity \\
  transport' = transport \\
  hasSubs' = hasSubs \\
  hasWriter' = hasWriter \\
  inboxNonEmpty' = inboxNonEmpty \\
  cascadeFor' = cascadeFor \\
  owner' = owner
\end{schema}

One change: the same \texttt{$cascadeFor = \emptyset$} guard \texttt{Connect}
gains, matching \texttt{identify\_client} joining \texttt{register\_client}
under \texttt{StoreLock.write()} in the design document's Section~4 fix.
\texttt{Renew} touches none of the three cascade-leg flags -- an explicit
\texttt{identify} call is contact, not \texttt{ensure\_writer}, and does not
by itself bind a writer.

\subsection{Activate --- the transport-owned side effects, one abstract step}
\label{sec:activate}

New in round 3. Models, collapsed into a single step, whatever combination
of a topic subscription, a writer binding, and an inbox write an ordinary
connected client performs after connecting -- \texttt{SubscriptionRegistry.subscribe},
\texttt{Hub.register\_writer}, and a peer's \texttt{publish} landing a
message in \texttt{c?}'s own inbox queue. The design document frames this
collapse explicitly ("modeled abstractly") for its cascade-incomplete
trace (Section~\ref{sec:fidelity-cascade-incomplete} below): I7 needs only
that these three facts can become true together, never which topic or
which message, so one schema setting all three suffices and no topic or
message carrier is introduced.

\begin{schema}{Activate}
  \Delta ConnReg \\
  c? : CONN
\where
  c? \in registered \\
  transport~c? = tup \\
  hasSubs' = hasSubs \oplus \{ c? \mapsto fset \} \\
  hasWriter' = hasWriter \oplus \{ c? \mapsto fset \} \\
  inboxNonEmpty' = inboxNonEmpty \oplus \{ c? \mapsto fset \} \\
  registered' = registered \\
  identity' = identity \\
  transport' = transport \\
  lease' = lease \\
  owner' = owner \\
  cascadeFor' = cascadeFor
\end{schema}

\texttt{Activate} does not renew the lease -- it is not named among the
contact-renews-lease operations in round 2's ruling, and keeping it inert on
\texttt{lease} is what keeps \texttt{LeaseLapse} reachable immediately
afterward in Section~\ref{sec:fidelity-cascade-incomplete}'s trace, exactly
as the design document's own sketch lists it. Nor does it check or touch
\texttt{cascadeFor}: while \texttt{c?} is mid its own
\texttt{TimedReapBegin}/\texttt{TimedReapEnd} window it has already left
\texttt{registered}, so \texttt{Activate}'s own guard already excludes it,
and no other connection's window has any bearing on what \texttt{c?} can do.

\subsection{TransportDies --- the environment, not the Hub}

Unchanged from round 1.

\begin{schema}{TransportDies}
  \Delta ConnReg \\
  c? : CONN
\where
  c? \in registered \\
  transport~c? = tup \\
  transport' = transport \oplus \{ c? \mapsto tdown \} \\
  registered' = registered \\
  identity' = identity \\
  lease' = lease \\
  owner' = owner \\
  hasSubs' = hasSubs \\
  hasWriter' = hasWriter \\
  inboxNonEmpty' = inboxNonEmpty \\
  cascadeFor' = cascadeFor
\end{schema}

\subsection{LeaseLapse --- time passes with nothing to renew it}
\label{sec:leaselapse}

\textbf{This is the one carry-over operation that changes.} Round 1
guarded \texttt{LeaseLapse} on \texttt{transport~c? = tdown}, reasoning
that a live transport implies contact keeps renewing the lease. H3 is the
proof that reasoning was wrong: \texttt{apply} is contact over a live
transport that, in the shipped code, renews nothing. The guard is
therefore weakened to drop the transport condition entirely -- a lease
lapses whenever its length has passed since the last renewal, live
transport or not, exactly matching \texttt{SessionLease.is\_live}, which
is pure time arithmetic and never inspects transport state.

\begin{schema}{LeaseLapse}
  \Delta ConnReg \\
  c? : CONN
\where
  c? \in registered \\
  lease~c? = llive \\
  lease' = lease \oplus \{ c? \mapsto llapsed \} \\
  registered' = registered \\
  identity' = identity \\
  transport' = transport \\
  owner' = owner \\
  hasSubs' = hasSubs \\
  hasWriter' = hasWriter \\
  inboxNonEmpty' = inboxNonEmpty \\
  cascadeFor' = cascadeFor
\end{schema}

\subsection{Read --- the non-destructive query}
\label{sec:read}

Models \texttt{HubClientRegistry.named\_sessions},
\texttt{live\_sessions}, and \texttt{repos} \emph{as this fix requires
them to behave}: a read changes nothing. \texttt{$\Xi$ ConnReg} is the
whole schema -- there is no parameter and no predicate beyond the
implicit \texttt{registered' = registered} (and every other component
held) that \texttt{$\Xi$} carries. This is I5's second half stated as a
structural fact about \texttt{Read}'s own definition, the same way I4 was
argued structurally in round~1: no reachability check is needed to show a
schema with no successor-changing predicate cannot change
\texttt{registered}. Section~\ref{sec:fidelity-h1}'s control is what
\texttt{Read} looks like without this fix.

\begin{schema}{Read}
  \Xi ConnReg
\end{schema}

\subsection{Depart --- the single removal-and-release coordinator}
\label{sec:depart}

Models \texttt{HubDisplay.drop\_connection}, fixed to do what its own
docstring today says it deliberately does not: release. One schema
stands for both triggers the ruling names as ``graceful'' and
``idle-reap,'' because both concretely call the identical
\texttt{drop\_connection} today, via
\texttt{session\_cleanup.py}'s teardown leg and
\texttt{lifecycle.disconnect\_connection} -- there is exactly one code
path, so there is exactly one schema. The guard is only
\texttt{$c? \in registered$}: no lease condition, because a session ending
cleanly or being idle-reaped by the SDK happens regardless of whether the
Hub's own lease has separately lapsed it. Removal and release are one
predicate, one step: there is no state in which \texttt{c?} has left
\texttt{registered} under this operation without \texttt{owner} already
reflecting the release.

\begin{schema}{Depart}
  \Delta ConnReg \\
  c? : CONN
\where
  c? \in registered \\
  registered' = registered \setminus \{ c? \} \\
  owner' = owner \oplus \\
  \quad~ (\{ s : SCENE | owner~s = conn(c?) \} \cross \{ unowned \}) \\
  hasSubs' = hasSubs \oplus \{ c? \mapsto fclear \} \\
  hasWriter' = hasWriter \oplus \{ c? \mapsto fclear \} \\
  inboxNonEmpty' = inboxNonEmpty \oplus \{ c? \mapsto fclear \} \\
  identity' = identity \\
  transport' = transport \\
  lease' = lease \\
  cascadeFor' = cascadeFor
\end{schema}

The three new lines are round 3's addition, in exactly the predicate that
already zeroes \texttt{owner} -- one atomic step, all four cascade legs at
once, matching \texttt{HubDisplay.drop\_connection} once it is fixed to run
the coordinator's full cascade (design document Section~2) rather than only
its first two legs. \texttt{Depart} needs no split: nothing in the design
document's Section~4 finding implicates it, and Section~\ref{sec:inv}
argues why its round-2 atomicity already forecloses the hazard that
motivates \texttt{TimedReap}'s split below.

\subsection{TimedReap --- the lease as a real background backstop, now in
two steps}
\label{sec:timedreap}

\texttt{Depart} and \texttt{Claim}'s embedded sweep are both
\emph{pull-based}: each fires only because some party -- a disconnecting
client, the SDK's idle-reap, or a peer's write -- happened to act.
Neither fires on its own. The bead this model exists to close reports a
connection that persisted 25.79 hours: no peer ever wrote again, and
nothing ever disconnected it, so neither pull-based path ever ran, and a
purely pull-based fix would have reproduced exactly that silence.
\texttt{TimedReap} realizes the one condition a background timer can act on
without any other party's help: the lease has lapsed. No other connection
needs to exist, act, or have ever existed, for it to become and stay
enabled.

Through round 2 this was one atomic schema, line-for-line \texttt{Depart}'s
with one added guard conjunct. Round 3 splits it into two:
\texttt{TimedReapBegin} performs exactly what round 2's \texttt{TimedReap}
did -- the registry-and-ownership removal, unchanged in shape from
\texttt{Depart}'s own first two legs -- and marks \texttt{c?} as
mid-cascade; \texttt{TimedReapEnd} performs the two new legs this round
adds (subscriptions and writer, inbox) and clears the mark. This is not a
cosmetic refactor. It is required to make the register\_client-vs-reap
race the design document reports (its Section~4) reachable by \fuzz\ 's
model at all -- the introduction above already argues why no ordering of
two fully atomic, freshly-guard-checked operations could ever witness it.
\texttt{Depart} and \texttt{Claim} are not split, because nothing reported
implicates them, and Section~\ref{sec:inv} argues their existing atomicity
already forecloses the same hazard without needing the device.

\begin{schema}{TimedReapBegin}
  \Delta ConnReg \\
  c? : CONN
\where
  c? \in registered \\
  lease~c? = llapsed \\
  cascadeFor = \emptyset \\
  registered' = registered \setminus \{ c? \} \\
  owner' = owner \oplus \\
  \quad~ (\{ s : SCENE | owner~s = conn(c?) \} \cross \{ unowned \}) \\
  cascadeFor' = \{ c? \} \\
  identity' = identity \\
  transport' = transport \\
  lease' = lease \\
  hasSubs' = hasSubs \\
  hasWriter' = hasWriter \\
  inboxNonEmpty' = inboxNonEmpty
\end{schema}

The guard's added conjunct, \texttt{$cascadeFor = \emptyset$}, is
\texttt{TimedReapBegin}'s own self-exclusion: only one connection's cascade
may be in flight at a time, matching \texttt{StoreLock}'s single-holder
property (\texttt{ConnReg}'s own \texttt{$\#cascadeFor \leq 1$} axiom
states the invariant this maintains). \texttt{hasSubs}, \texttt{hasWriter},
and \texttt{inboxNonEmpty} are deliberately left untouched here -- this is
the registry-and-ownership removal alone, matching real
\texttt{\_depart\_lapsed}'s own first act (design document Section~1: ``C is
popped from \texttt{\_sessions}. \texttt{StoreLock} is still held.'') --
the window this exposes, and what closes it, is the entire point of the
split.

\begin{schema}{TimedReapEnd}
  \Delta ConnReg \\
  c? : CONN
\where
  cascadeFor = \{ c? \} \\
  cascadeFor' = \emptyset \\
  hasSubs' = hasSubs \oplus \{ c? \mapsto fclear \} \\
  hasWriter' = hasWriter \oplus \{ c? \mapsto fclear \} \\
  inboxNonEmpty' = inboxNonEmpty \oplus \{ c? \mapsto fclear \} \\
  registered' = registered \\
  identity' = identity \\
  transport' = transport \\
  lease' = lease \\
  owner' = owner
\end{schema}

\texttt{TimedReapEnd} models the cascade tail this round adds --
\texttt{Hub.on\_disconnect} (subscriptions, writer) and the transport
sink (inbox) -- and is enabled precisely when some connection's cascade is
open, closing it unconditionally once fired: no other precondition, so it
cannot be indefinitely deferred once \texttt{TimedReapBegin} has run
(Section~\ref{sec:verify}'s deadlock re-run confirms this rather than
assuming it). Together, \texttt{TimedReapBegin}\,;\,\texttt{TimedReapEnd} in
immediate succession, with nothing else interleaved, compute exactly what
round 2's single \texttt{TimedReap} computed, plus the three new
\(\oplus\) terms -- the split changes nothing about the property round 2
already proved; it only creates the one window this round needs to state
and then close a new one.

\subsection{Claim --- the ownership-gated write, renewing itself first}
\label{sec:claim}

Models \texttt{HubDisplay.apply}: \texttt{SetProperty}/\texttt{RemoveElement}'s
\texttt{OwnerTracker.require\_ownership} check, and (folded into the same
guard, as in round 1) \texttt{AddElement}'s grant onto a not-yet-owned
scene. \textbf{Three things happen in one atomic step, in this order,
exactly matching \texttt{apply}'s own docstring:}

\begin{enumerate}
  \item \textbf{Renew the caller.} \texttt{$lease' = lease \oplus
    \{c? \mapsto llive\}$} happens unconditionally, before anything else is
    computed. This is the ruling's third decision -- an authenticated
    write is contact, and contact renews -- and it is the one line the
    shipped code is missing (\texttt{apply} never calls
    \texttt{HubClientRegistry.record} for its own caller).
  \item \textbf{Sweep every \emph{other} lapsed connection.} The set
    \(\{c2 : CONN \mathbin{|} c2 \in registered \land c2 \neq c? \land
    lease~c2 = llapsed\}\) names every connection but the caller whose
    lease has lapsed; each is dropped from \texttt{registered} and every
    scene it owned is released to \texttt{unowned}, in the same step as
    the renewal above. The \(c2 \neq c?\) exclusion is what step 1 buys:
    because \texttt{c?} was just renewed to \texttt{llive}, it could never
    satisfy this set's own guard anyway, but stating the exclusion
    explicitly is what distinguishes this schema from
    Section~\ref{sec:fidelity-h3}'s control, where the exclusion is
    exactly what is missing.
  \item \textbf{Claim, against the swept result.} The ownership guard --
    \texttt{unowned} or already the caller's -- is evaluated against
    \texttt{owner} \emph{after} step 2's release, matching
    \texttt{apply}'s docstring: ``reaps and releases first, before the
    ownership check.'' A reconnecting client claiming a scene a departed,
    same-call-swept predecessor held is not shadowed by it.
\end{enumerate}

\texttt{transport~c? = tup} is required as a precondition, matching
\texttt{Renew}: an authenticated write, like any authenticated contact,
can only arrive over a live transport.

\begin{schema}{Claim}
  \Delta ConnReg \\
  c? : CONN \\
  s? : SCENE
\where
  c? \in registered \\
  transport~c? = tup \\
  lease' = lease \oplus \{ c? \mapsto llive \} \\
  registered' = registered \setminus \\
  \quad~ \{ c2 : CONN | c2 \in registered \land c2 \neq c? \land
    lease~c2 = llapsed \} \\
  (owner \oplus (\{ s : SCENE | \exists c2 : CONN @ \\
    \qquad c2 \in registered \land c2 \neq c? \land lease~c2 = llapsed
    \land owner~s = conn(c2) \} \\
    \qquad \cross \{ unowned \}))~s? = unowned \\
  \lor \\
  (owner \oplus (\{ s : SCENE | \exists c2 : CONN @ \\
    \qquad c2 \in registered \land c2 \neq c? \land lease~c2 = llapsed
    \land owner~s = conn(c2) \} \\
    \qquad \cross \{ unowned \}))~s? = conn(c?) \\
  owner' = (owner \oplus (\{ s : SCENE | \exists c2 : CONN @ \\
    \qquad c2 \in registered \land c2 \neq c? \land lease~c2 = llapsed
    \land owner~s = conn(c2) \} \\
    \qquad \cross \{ unowned \})) \oplus \{ s? \mapsto conn(c?) \} \\
  hasSubs' = hasSubs \oplus (\{ c2 : CONN | c2 \in registered \land
    c2 \neq c? \land lease~c2 = llapsed \} \cross \{ fclear \}) \\
  hasWriter' = hasWriter \oplus (\{ c2 : CONN | c2 \in registered \land
    c2 \neq c? \land lease~c2 = llapsed \} \cross \{ fclear \}) \\
  inboxNonEmpty' = inboxNonEmpty \oplus (\{ c2 : CONN | c2 \in registered
    \land c2 \neq c? \land lease~c2 = llapsed \} \cross \{ fclear \}) \\
  identity' = identity \\
  transport' = transport \\
  cascadeFor' = cascadeFor
\end{schema}

The three new lines mirror the \(\oplus\) already computing \texttt{owner'}
exactly: same sweep set, same \(c2 \neq c?\) exclusion (the exclusion that
closes H3/I6 is untouched -- the caller is never among the connections
its own write clears), applied to the three new flags instead of ownership.
\texttt{Claim} needs no \texttt{cascadeFor} guard of its own: its own guard
already requires \texttt{$c? \in registered$}, which a connection currently
mid \texttt{TimedReapBegin}/\texttt{TimedReapEnd} has already lost, so
\texttt{Claim(c?, s?)} for that \texttt{c?} is disabled by its existing
guard with no new conjunct needed, and \texttt{Claim} for any \emph{other}
connection has no state in common with a cascade window that names a
different \texttt{c?} -- adding the guard there would exclude no
additional trace.

% ============================================================
\section{Invariants}
\label{sec:inv}
% ============================================================

None of the six properties is placed in the \texttt{ConnReg} schema, for
the same reason round 1 gave: a predicate there would be enforced as a
guard on every successor, silently disabling the very operation that
would break it.

\paragraph{I1 --- transport-gone implies \emph{unconditionally} eventually
reaped.}
As in round 1, this is a liveness property argued over enabledness and
reachability, not a single-state safety formula -- Z has no temporal
operator of its own, so the claim is: for any \texttt{c?} with
\texttt{transport~c? = tdown}, no reachable state traps \texttt{c?} in
\texttt{registered} forever, and once it leaves, every scene it owned is
already released. \texttt{Renew}'s guard
means a connection can never re-arm its lease once its transport has died
without an intervening \texttt{Connect}. Once
\texttt{$transport~c? = tdown \land lease~c? = llive$}, \texttt{LeaseLapse}
is enabled and depends on nothing but \texttt{c?}'s own state; once
\texttt{lease~c? = llapsed}, \texttt{TimedReapBegin} (Section
\ref{sec:timedreap}) is enabled and \emph{also} depends on nothing but
\texttt{c?}'s own state and \texttt{cascadeFor} being empty -- an
uncontested condition for a lone connection, since nothing else can be
mid-cascade in a carrier holding only \texttt{c?} -- not on any other
connection existing, not on any read ever running, not on any disconnect
ever being detected. This is
the strengthening this addendum makes, and it is worth being exact about
what was actually wrong with the previous round's argument, because it is
not a reachability gap. \texttt{Depart}'s guard, \texttt{$c? \in
registered$}, was \emph{already} satisfiable with no second connection and
no \texttt{Read} in the picture -- a \texttt{fuzz} reachability check
alone cannot distinguish ``this operation's formal precondition holds''
from ``something in the real system will actually invoke it,'' because Z
enabledness is silent on triggers. The error in trusting that chain as
\emph{unconditional} was informal, not formal: this document's own prose
already restricted \texttt{Depart} to modelling a real disconnect signal
or the SDK's idle-reap (Section~\ref{sec:depart}), and for a killed
process neither ever arrives -- so \texttt{Depart}'s formal availability
was never in question, only whether anything in the real system would
ever exercise it. \texttt{TimedReap} closes that informal gap by being an
operation whose intended trigger -- time passing -- genuinely requires
nothing external, which is a claim about what invokes it in the real
\texttt{HubClientRegistry}, not one \texttt{fuzz} or \texttt{probcli}
can check for either operation. What \emph{is} newly, formally checkable
is that this reading is now backed by an operation whose guard depends
on nothing but \texttt{c?}'s own recorded state, so a periodic caller
with no reference to any other connection's identity, and no external
signal delivered to it at all, could implement exactly this guard: the
chain
\texttt{TransportDies}\,;\,\texttt{LeaseLapse}\,;\,\texttt{TimedReapBegin}\,;
\,\texttt{TimedReapEnd},
run against a carrier holding only \texttt{c?} itself -- no second
connection, no \texttt{Read}, no \texttt{Depart} -- is a complete,
self-contained witness that \texttt{c?} leaves \texttt{registered} with
every scene it owned released and every cascade-leg flag cleared, and
Section~\ref{sec:verify} confirms this
trace by reachability with \texttt{CONN} instantiated to a single element.
\texttt{Depart} and \texttt{Claim}'s embedded sweep remain additional,
faster dischargers of a lapsed lease when their own triggers happen to
fire first (Section~\ref{sec:serialize} shows the three do not race); they
are no longer I1's \emph{only} route, which is what round 1's original
coverage audit named as an open gap and this addendum closes.

\paragraph{I2 --- reap releases ownership.}
\[
  \lnot\, \exists\, s : SCENE; c : CONN @
    owner~s = conn(c) \land c \notin registered.
\]
Restated verbatim from round 1. It remains true, but it is no longer the
whole of the argument, because there are now three operations that can
remove a connection from \texttt{registered}
(\texttt{Depart}, \texttt{TimedReapBegin}, and \texttt{Claim}'s embedded
sweep of others) instead of one, plus one that must never remove anyone
(\texttt{Read}). I5, below, is the version of this argument that
actually accounts for all four.

\paragraph{I3 --- at most one live owner per identity.}
\[
  \lnot\, \exists\, s : SCENE; c, c' : CONN @ \\
  \quad~ c \in registered \land c' \notin registered \land \\
  \quad~ c' \in \dom identity \land identity~c = identity~c' \land \\
  \quad~ c \neq c' \land owner~s = conn(c').
\]
Unchanged in statement from round 1: I5's general form already rules out
\emph{any} scene being owned by a connection outside \texttt{registered},
so I3 holds as the same-identity corollary of I5 that round 1 argued it
was of I2.

\paragraph{I4 --- reconnect inherits nothing special (settled).}
Unchanged from round 1 -- \texttt{Connect}'s definition has not changed,
and the argument is unchanged: no successor of \texttt{Connect} touches
any other connection's state, so a reconnecting client's only route to
ownership is the same \texttt{Claim} every connection uses.

\paragraph{I5 --- leaving the registry iff released, however triggered.}
\[
  \lnot\, \exists\, s : SCENE; c : CONN @
    owner~s = conn(c) \land c \notin registered
\]
\emph{together with:} no operation whose only effect is a read (that is,
\texttt{Read} itself) ever changes \texttt{registered}.

The first conjunct is I2's formula, now proven against four removal
paths instead of one: \texttt{Depart} and \texttt{TimedReapBegin} each
release every scene \texttt{c?} owned in the same step they remove
\texttt{c?} (Sections~\ref{sec:depart} and~\ref{sec:timedreap} -- the two
predicates are identical, registry-and-ownership-wise, apart from
\texttt{TimedReapBegin}'s added lease guard and its
\texttt{$cascadeFor = \emptyset$} self-exclusion, so this argument covers
both at once; the split's second half, \texttt{TimedReapEnd}, touches
neither \texttt{registered} nor \texttt{owner} at all, so it plays no part
in I2 or I5), and \texttt{Claim}'s embedded sweep
releases every scene each swept \texttt{c2} owned in the same step it
removes \texttt{c2} from \texttt{registered} (Section~\ref{sec:claim}'s
step 2). No other operation ever assigns \texttt{owner} to a connection
outside \texttt{registered} -- \texttt{Claim}'s own guard,
\texttt{$c? \in registered$}, and its postcondition,
\texttt{$c? \in registered'$} always (Section~\ref{sec:claim}'s
\(c2 \neq c?\) exclusion guarantees the caller is never among the
connections it removes), rule that out at the point of assignment. The
second conjunct is \texttt{Read}'s own definition
(Section~\ref{sec:read}): \texttt{$\Xi$ ConnReg} carries
\texttt{registered' = registered} by construction, so this half needs no
reachability check, exactly as I4 needed none.

This is what closes H1 and H2 together, as one invariant rather than two
patches: H1 broke the second conjunct (a read removed someone), H2 broke
the first conjunct via a specific operation (\texttt{Depart} removing
without releasing); I5 states the shape both holes share -- \emph{leaving
the registry and releasing ownership are the same event, no matter which
of the model's removing operations produced it, and a non-removing
operation produces neither} -- so that a fourth departure trigger
appearing later in the code would have to be checked against the same
general property, not against a fourth ad hoc patch. \texttt{TimedReap}
\emph{is} that fourth trigger, added within this same round, and I5's
statement did not need to change at all to absorb it -- only its proof
gained one more, textually near-identical, case.

\paragraph{I6 --- a live writer keeps its ownership.}
\[
  \lnot\, \exists\, s : SCENE; c : CONN @
    owner~s = conn(c) \land c \notin registered \land transport~c = tup.
\]
I5's negation, specialized to the case where the departing connection's
transport never died. Every witness to I5's general formula that
\texttt{Depart} could produce requires nothing of transport state at all
(Section~\ref{sec:depart}'s guard is orthogonal to it), so it does not
distinguish this case on its own -- but that is exactly why \texttt{c?}'s
own \texttt{transport} being \texttt{tup} throughout a trace is a
meaningful discriminator elsewhere: a \texttt{TimedReapBegin}-produced
witness \emph{also} requires nothing of transport state, so a witness
confined to \texttt{transport~c? = tup} could still, in principle, come
from either \texttt{Depart} or \texttt{TimedReapBegin} acting on a
connection that never crashed at all -- an ordinary graceful departure or
a lease that genuinely lapsed from real inactivity, neither of which is a
bug. What I6 rules out is narrower and sharper: that \texttt{Claim}'s own
sweep, not \texttt{Depart} or \texttt{TimedReapBegin}, ever catches its
\emph{own caller} on the very call that caller is making -- the one
mechanism H3 names, and the one no legitimate departure operation is even
attempting to perform mid-write. In the fixed spec that cannot happen,
because \texttt{Claim}'s \(c2 \neq c?\) exclusion (Section~\ref{sec:claim},
step 2) means the caller can never appear in the set it sweeps, regardless
of what its own \texttt{lease} showed on entry. Section~\ref{sec:fidelity-h3}'s
control removes exactly that exclusion and shows the formula above becomes
reachable in three steps that never call \texttt{TransportDies},
\texttt{Depart}, \texttt{TimedReapBegin}, or \texttt{TimedReapEnd} at all --
isolating the claim to \texttt{Claim} alone, exactly as the
fidelity-control discipline requires.

\paragraph{I7 --- departure clears the full cascade, however triggered.}
\[
  \lnot\, \exists\, c : CONN @ \\
  \quad~ c \notin registered \land c \notin cascadeFor \land \\
  \quad~ (hasSubs~c = fset \lor hasWriter~c = fset \lor
    inboxNonEmpty~c = fset).
\]
The general form of I5, extended from the registry-and-ownership leg to
all three of the legs round 3 adds. The \texttt{$c \notin cascadeFor$}
conjunct is the one genuinely new piece of the statement, and it earns a
careful defense rather than a wave: without it, I7 would be reported
\texttt{FOUND}-violated even in this, the fixed, spec, purely as an
artifact of the \texttt{TimedReapBegin}/\texttt{TimedReapEnd} split --
immediately after \texttt{TimedReapBegin(c?)}, \texttt{c?} has already
left \texttt{registered} but \texttt{hasWriter(c?)} may still read
\texttt{fset}, because clearing it is \texttt{TimedReapEnd}'s job, not
\texttt{TimedReapBegin}'s. That is not a defect in the fix; it is the
exposed half of the very window this round's split exists to reveal, and
it is invisible to every external operation for the same reason round 2's
``two-lock concern is moot by construction'' argument already gives for
\emph{every} atomic schema in this document: nothing else can act on
\texttt{c?} while its own cascade is open (\texttt{Depart}'s and
\texttt{Claim}'s guards on \texttt{c?} already fail the moment
\texttt{TimedReapBegin} has run, and \texttt{Connect}/\texttt{Renew}'s new
\texttt{$cascadeFor = \emptyset$} guard, argued below, forecloses the one
operation that otherwise could). I7 is stated, correctly, over states an
external observer can actually reach between complete operations -- exactly
where I2, I5, and I6 are already implicitly stated, since Z's own
atomicity means none of them has ever needed to say ``except mid-Depart''
or ``except mid-Claim.'' \texttt{TimedReapBegin}/\texttt{TimedReapEnd} is
the one place this round deliberately widens the grain past a single
schema, so it is the one place the exemption must be written down rather
than left implicit. The exemption is vacuous for every connection that
ever departs via \texttt{Depart} or \texttt{Claim}'s sweep, since neither
operation ever adds anything to \texttt{cascadeFor} -- \texttt{$c
\notin cascadeFor$} is unconditionally true for them, so I7 binds on them
with no relaxation at all. This is what closes the gap Section~1 of the
design document names: \texttt{Depart} already reached all four cascade
legs once fixed; \texttt{TimedReapBegin}/\texttt{TimedReapEnd} and
\texttt{Claim}'s sweep now do too.

\paragraph{The two-lock concern is moot by construction.}
The current, unfixed code holds \texttt{HubClientRegistry}'s own lock
around registration and \texttt{OwnerTracker} under no lock of its own
(\texttt{HubDisplay}'s store lock covers it instead) -- two different
locks, held by two different objects, and H1/H2 are exactly what happens
when a caller updates one without the other under one critical section.
This model does not need a second lock in its own terms, because it does
not need coordinating in the first place: \texttt{Read} never mutates
\texttt{registered} (Section~\ref{sec:read}), so the only three
operations that ever do -- \texttt{Depart}, \texttt{TimedReapBegin}, and
\texttt{Claim} -- are also, by their own single predicate, the only
three operations that ever mutate \texttt{owner} on a departure. There
is no interleaving in which one without the other could be observed,
because in this model, as the fix requires of the real
\texttt{HubDisplay} (the one object holding both \texttt{\_clients} and
\texttt{\_owners}), removing-from-the-registry and releasing-ownership
are not two updates that a lock discipline keeps in step; they are one
update.

\paragraph{Bringing \texttt{Connect}/\texttt{Renew} under the same
discipline introduces no new lock-acquisition cycle.}
\label{sec:storelock-nesting}
The design document's Section~4 fix nests two real locks on a path that
previously took only the inner one: \texttt{register\_client}/
\texttt{identify\_client} currently take only \texttt{HubClientRegistry}'s
own internal lock; the fix additionally wraps them in
\texttt{HubDisplay.\_lock.write()} (\texttt{StoreLock}), so the call now
holds \texttt{StoreLock} \emph{and then} the inner registry lock, nested,
exactly as \texttt{apply}/\texttt{replace\_scene}/\texttt{show\_scene}
already do. A locking-discipline change that nests two locks is exactly
the class of change this project's z-spec mandate requires proving
deadlock-free, not assuming it. The argument is short because the nesting
order does not change: on every path, before and after this fix,
\texttt{StoreLock} is acquired first and the inner \texttt{HubClientRegistry}
lock (fully absorbed into this model's schema atomicity, per the paragraph
above) second, never the reverse -- \texttt{register\_client}/
\texttt{identify\_client} joining that discipline is one more caller taking
the locks in the \emph{same} order every existing caller already does, not
a new order for anyone to take them in. A deadlock needs a cycle in the
lock-acquisition graph; adding an edge in a direction every existing edge
already points cannot create one. In this model, that global order is
represented by \texttt{cascadeFor}: it is acquired by exactly one operation
(\texttt{TimedReapBegin}), released by exactly one operation
(\texttt{TimedReapEnd}), and merely tested, never acquired, by
\texttt{Connect} and \texttt{Renew} -- a single-owner, single-releaser
resource is structurally incapable of participating in a cycle with
itself, since a cycle needs at least two resources each waiting on the
other. Section~\ref{sec:verify}'s full \texttt{-model\_check}, re-run
against this round's extended carrier, confirms no deadlock rather than
resting on this structural argument alone.

\paragraph{The three coordinators serialize; none of them races.}
\label{sec:serialize}
Round 2's implementation-facing question is whether a background timer
firing \texttt{TimedReap} can race \texttt{Depart} (an external disconnect
arriving at the same moment) or \texttt{Claim} (a write arriving at the
same moment) on the \emph{same} connection, and land the registry or the
ownership table in a state neither operation intended. In this model the
question has a structural answer, not merely an empirical one: a Z
operation schema denotes a single, indivisible relation between a
before-state and an after-state, and \texttt{DEFAULT\_SETSIZE}-bounded
model-checking with ProB explores every possible \emph{interleaving} of
enabled operations, never a partial or torn application of one. There is
no notion in this model of two operations executing ``at once'' and
observing each other mid-update -- only of one full step happening, then
the next full step choosing among whatever is enabled in the resulting
state. Three consequences follow, and Section~\ref{sec:verify}'s full
\texttt{-model\_check} (which visits every reachable state and reports no
deadlock) is what confirms none of them is merely assumed:

\begin{enumerate}
  \item \textbf{Whichever fires first wins, and correctly.} If
    \texttt{TimedReapBegin(c?)} and \texttt{Depart(c?)} are \emph{both}
    enabled in some state (a lapsed lease and a live,
    graceful-disconnect-eligible connection are not mutually exclusive
    conditions), firing either one removes \texttt{c?} from
    \texttt{registered} and releases every scene it owned -- the two
    predicates differ only in \texttt{TimedReapBegin}'s extra guard, never
    in their registry-and-ownership effect. The second operation, whichever
    it is, is then simply disabled: its guard, \texttt{$c? \in registered$},
    fails, because the first one already removed \texttt{c?}. There is no
    state in which both fire on the same \texttt{c?} and no state in
    which their effects could conflict, because they compute the
    identical postcondition on \texttt{registered} and \texttt{owner}.
  \item \textbf{A write is never stranded behind a timer, or vice versa.}
    \texttt{Claim(c2, s?)} for a \emph{different} connection \texttt{c2}
    and \texttt{TimedReapBegin(c1)} for a lapsed \texttt{c1} are
    independent: neither's guard reads the other's parameter, and
    \texttt{Claim}'s own embedded sweep (Section~\ref{sec:claim}, step 2)
    would perform the identical release on \texttt{c1} if
    \texttt{TimedReapBegin} had not already done so first. Whichever fires
    first, the reachable state after both have run (in either order) is
    the same: \texttt{c1} departed and released, \texttt{c2}'s claim
    resolved against a registry that no longer shows \texttt{c1} as a
    candidate to be swept. Reachability is monotonic under adding
    \texttt{TimedReap} to a spec that already had \texttt{Depart} and
    \texttt{Claim}: every state and trace reachable before this addendum
    remains reachable after it, because no existing operation's guard or
    effect changed.
  \item \textbf{A connection can never claim while being timed-reaped
    out from under itself.} \texttt{Claim(c?, s?)} renews \texttt{c?}'s
    own lease to \texttt{llive} as its first effect
    (Section~\ref{sec:claim}, step 1), and \texttt{TimedReapBegin(c?)}
    requires \texttt{lease~c? = llapsed}. Because both are whole,
    indivisible steps, there is no state in which \texttt{Claim(c?, s?)}
    is ``in progress'' for \texttt{TimedReapBegin(c?)} to observe a stale
    lease during: either \texttt{Claim} has already completed, in which
    case \texttt{c?}'s lease is \texttt{llive} and \texttt{TimedReapBegin(c?)}
    is disabled, or it has not yet run at all.
\end{enumerate}

None of this needs a fourth invariant for \texttt{Depart}/\texttt{Claim}:
it is I5 and I6, reread as claims about \emph{any} pair of those two
registry-mutating operations plus \texttt{TimedReapBegin} rather than
about \texttt{Claim} alone, and all three are already checked exhaustively
over every reachable state, from every reachable ordering, by
Section~\ref{sec:verify}'s I5/I6 goal checks. A genuine race among
\emph{these three} producing an inconsistent intermediate state would
surface there directly, as a \texttt{FOUND} counterexample naming the
offending trace -- not merely as a deadlock, which
Section~\ref{sec:verify}'s separate full model-check also confirms absent,
so neither failure mode is left unchecked for this trio.

\paragraph{\texttt{Connect}/\texttt{Renew} are the one pair this argument
does \emph{not} cover, and closing them needs the new guard, not just the
atomicity argument above.}
\label{sec:reconnect-race}
The three consequences above rest on every participant being a single,
whole, indivisible Z step -- which is exactly what makes
\texttt{TimedReapBegin}/\texttt{TimedReapEnd}'s deliberate split the one
place in this document where that rest is \emph{not} available for free.
\texttt{Connect(c?, i?)} sets \texttt{$hasWriter'(c?) = fset$}
unconditionally (Section~\ref{sec:depart}'s sibling, \texttt{Connect}'s own
definition above); nothing else in the fixed spec ever sets
\texttt{hasWriter} back to \texttt{fclear} for a connection that remains in
\texttt{registered} -- the only two operations that ever do,
\texttt{TimedReapEnd} and \texttt{Depart}, either presuppose or produce
\texttt{c?} being outside \texttt{registered} in the very same or a prior
step. So the state
\texttt{$c? \in registered \land hasWriter(c?) = fclear$} is reachable in
the fixed spec if, and only if, \texttt{TimedReapEnd(c?)} is permitted to
fire \emph{after} a \texttt{Connect(c?, \ldots)} that reconnected \texttt{c?}
during \texttt{c?}'s own still-open cascade window -- precisely the
interleaving \texttt{cascadeFor} exists to name and the new
\texttt{$cascadeFor = \emptyset$} guard on \texttt{Connect}/\texttt{Renew}
exists to forbid: once \texttt{TimedReapBegin(c?)} has set
\texttt{$cascadeFor = \{c?\}$}, \texttt{Connect}/\texttt{Renew} for
\emph{any} connection are disabled until \texttt{TimedReapEnd} clears it,
so no reconnect can land inside the window at all, and by the time one
can, \texttt{hasWriter(c?)} is already settled at whatever
\texttt{TimedReapEnd} left it (\texttt{fclear}) with nothing left in
flight to overwrite. Section~\ref{sec:fidelity-reconnect-race}'s control
removes exactly this one guard and shows the goal above becomes reachable
in the shortest witness the design document's own trace names -- isolating
the hazard to the missing guard alone, exactly as the fidelity-control
discipline requires. This is the one place in round 3 where atomicity
alone, the argument this whole section otherwise rests on, is not
sufficient, and it is not sufficient for a reason worth stating plainly:
\texttt{Connect}/\texttt{Renew} are the real code's one gap relative to
its own stated discipline (design document Section~4, quoting
\texttt{hub\_display.py}'s docstring: ``every write runs under
\texttt{StoreLock}''), so they are the one pair this model must show a
race \emph{without} the fix and its absence \emph{with} it, rather than
argue away structurally as it does for every other pair.

% ============================================================
\section{Verification}
\label{sec:verify}
% ============================================================

Type-checking is a \texttt{fuzz} pass:
\begin{verbatim}
  fuzz docs/connection_lease_reaping.tex
\end{verbatim}

I2, I3, I5, and I6 are safety properties, checked by asking ProB whether
their negation is \emph{reachable}. A goal reported \emph{not found} means
the invariant holds on every reachable state.

\begin{verbatim}
  # I2 (goal NOT found in this, the fixed, spec)
  probcli docs/connection_lease_reaping.tex -model_check -nodead \
    -goal "#s.(#c.(s:SCENE & c:CONN & owner(s) = conn(c) &
                   not(c : registered)))" \
    -p DEFAULT_SETSIZE 3

  # I3 (goal NOT found in this, the fixed, spec)
  probcli docs/connection_lease_reaping.tex -model_check -nodead \
    -goal "#s.(#c.(#c2.(s:SCENE & c:CONN & c2:CONN &
                        c : registered & not(c2 : registered) &
                        c2 : dom(identity) & identity(c) = identity(c2) &
                        not(c = c2) & owner(s) = conn(c2))))" \
    -p DEFAULT_SETSIZE 3

  # I5 (identical goal to I2 -- I5's first conjunct; goal NOT found)
  probcli docs/connection_lease_reaping.tex -model_check -nodead \
    -goal "#s.(#c.(s:SCENE & c:CONN & owner(s) = conn(c) &
                   not(c : registered)))" \
    -p DEFAULT_SETSIZE 3

  # I6 (goal NOT found in this, the fixed, spec)
  probcli docs/connection_lease_reaping.tex -model_check -nodead \
    -goal "#s.(#c.(s:SCENE & c:CONN & owner(s) = conn(c) &
                   not(c : registered) & transport(c) = tup))" \
    -p DEFAULT_SETSIZE 3

  # I7 (goal NOT found in this, the fixed, spec)
  probcli docs/connection_lease_reaping.tex -model_check -nodead \
    -goal "#c.(c:CONN & not(c : registered) & not(c : cascadeFor) &
               (hasSubs(c) = fset or hasWriter(c) = fset or
                inboxNonEmpty(c) = fset))" \
    -p DEFAULT_SETSIZE 3
\end{verbatim}

Round 2 confirmed I2/I3/I5/I6 exhaustively at \texttt{DEFAULT\_SETSIZE 3}
(313{,}093 states) against round 2's own spec. That earlier count does not
automatically carry over to round 3's spec -- \texttt{Connect} and
\texttt{Renew} gained a new guard, \texttt{TimedReap} became two
operations, and every operation gained four new state components -- so the
same, structurally different model needs its own confirmation.
\textbf{This round's own exhaustive confirmation of I2, I3, I5, I6, and I7
is at \texttt{DEFAULT\_SETSIZE 2} (4{,}705 states, ten operations
covered, no deadlock).} \texttt{DEFAULT\_SETSIZE 3} was attempted for I7
and for the reconnect-race fidelity control's goal
(Section~\ref{sec:fidelity-reconnect-race}) and did not complete: the
enlarged carrier passed 380{,}000 states at 25--27\% coverage, with the
running total still climbing, before the search was stopped. This is
stated plainly, not glossed over: \texttt{DEFAULT\_SETSIZE 3} exhaustive
confirmation of I2/I3/I5/I6/I7 over round 3's carrier is an open
follow-up, and the commands above are given at \texttt{DEFAULT\_SETSIZE
3} as the target check for whenever that follow-up runs, not as a claim
already discharged this round.

I1's mechanism is confirmed by reachability of its positive outcome:

\begin{verbatim}
  # I1's chain completes (goal FOUND): a connection that connected, died,
  # and departed is outside registered with a dead transport
  probcli docs/connection_lease_reaping.tex -model_check -nodead \
    -goal "#c.(c:CONN & c : dom(identity) & not(c : registered) &
               transport(c) = tdown)" \
    -p DEFAULT_SETSIZE 3
\end{verbatim}

The state a witness ends in does not, by itself, say which operation
produced it: \texttt{Depart} and \texttt{TimedReapBegin} compute the
identical postcondition on \texttt{registered} and \texttt{owner}
(Section~\ref{sec:timedreap}), so a reachability goal alone cannot certify
that \texttt{TimedReap} specifically was exercised. The
same goal remains \emph{FOUND} at \texttt{DEFAULT\_SETSIZE 1}, which
bounds \texttt{CONN} to a single element and so rules out any second
connection appearing anywhere in the witness by construction of the
carrier:

\begin{verbatim}
  # goal FOUND even with only one connection in the whole carrier
  probcli docs/connection_lease_reaping.tex -model_check -nodead \
    -goal "#c.(c:CONN & c : dom(identity) & not(c : registered) &
               transport(c) = tdown)" \
    -p DEFAULT_SETSIZE 1
\end{verbatim}

The printed witness is not reproducible evidence on its own -- ProB's
search order is not obliged to return the same trace on every run, and a
run of this exact check returned a 4-step \texttt{Depart}-based trace on
one invocation and a 6-step one passing through \texttt{Claim} and
\texttt{LeaseLapse} on another, because \texttt{Depart}'s guard,
\texttt{$c? \in registered$} alone, was already satisfiable the moment
\texttt{TransportDies} fires and does not require the lease to have
lapsed at all. What is reproducible, and what actually proves the
strengthening, is the structural argument already given above: at any
point where \texttt{lease~c? = llapsed} holds for a registered \texttt{c?}
-- which \texttt{LeaseLapse}'s postcondition establishes unconditionally,
independent of any other connection -- \texttt{TimedReapBegin(c?)} is
enabled by inspection of its own guard (once \texttt{cascadeFor} is empty,
uncontested for a lone connection), whether or not the same state also
happens to enable \texttt{Depart}. A search tool that reports \emph{some}
witness cannot be asked to prefer one enabled operation over another
that produces the identical effect; the schema's own precondition is
what the proof rests on, not which branch a particular run's search
order happened to take.

The fix's positive outcome -- a live connection successfully claims a
scene its departed, same-identity predecessor used to own -- is also
confirmed reachable:

\begin{verbatim}
  # A live connection owns a scene a departed same-identity predecessor
  # used to own (goal FOUND).
  probcli docs/connection_lease_reaping.tex -model_check -nodead \
    -goal "#s.(#c.(#c2.(s:SCENE & c:CONN & c2:CONN & c : dom(identity) &
                        not(c = c2) & identity(c) = identity(c2) &
                        not(c : registered) & c2 : registered &
                        owner(s) = conn(c2))))" \
    -p DEFAULT_SETSIZE 3
\end{verbatim}

\paragraph{Deadlock.} Round 1/2's argument was that \texttt{Connect} carries
no precondition beyond its two arguments' types, so it is enabled in every
reachable state. Round 3 adds one: \texttt{Connect} now also requires
\texttt{$cascadeFor = \emptyset$}, so that argument alone no longer
suffices on its own -- it is exactly the nesting
Section~\ref{sec:storelock-nesting} argues introduces no new
lock-acquisition cycle, because \texttt{TimedReapEnd} (the sole releaser of
\texttt{cascadeFor}) carries no precondition beyond \texttt{$cascadeFor =
\{c?\}$} and is therefore always enabled the moment
\texttt{TimedReapBegin} has run, so \texttt{Connect} is never blocked
longer than one step. The full model-check, re-run against this round's
extended carrier (ten operations, four new state components), confirms no
deadlock at \texttt{DEFAULT\_SETSIZE 2} (4,705 states, all visited) rather
than resting on the structural argument alone:

\begin{verbatim}
  probcli docs/connection_lease_reaping.tex -model_check -p DEFAULT_SETSIZE 2
\end{verbatim}

\texttt{DEFAULT\_SETSIZE 3} was attempted and did not finish within this
round's session: the enlarged carrier -- three new total flag-functions
plus \texttt{cascadeFor} -- had already passed 380{,}000 states at
25--27\% coverage, well beyond round 2's own 313{,}093-state total at the
same setsize, with the running total still climbing when the search was
stopped. Stated plainly rather than glossed over: \texttt{DEFAULT\_SETSIZE
2}'s exhaustive result is what this round's deadlock claim rests on;
\texttt{DEFAULT\_SETSIZE 3} exhaustive confirmation is an open follow-up.

% ============================================================
\section{Fidelity: reproducing five shipped holes}
\label{sec:fidelity}
% ============================================================

A model that cannot reproduce a known defect proves nothing about closing
it. Each hole gets its own control, changing exactly one thing from this,
the fixed, spec -- so that a goal found in one control and not found in
the fixed spec attributes the defect to that one change and no other.
Round 2 closed three (H1--H3, below); round 3 adds two more --
Section~\ref{sec:fidelity-cascade-incomplete}'s cascade-incomplete control,
reproducing round 3's own newly-added gap, and
Section~\ref{sec:fidelity-reconnect-race}'s reconnect-races-reap control,
reproducing the register\_client-vs-reap hazard the design document's
Section~4 reports.

\subsection{H1 --- destructive-read-sweep}
\label{sec:fidelity-h1}

\texttt{docs/connection\_lease\_reaping\_destructive\_read\_buggy.tex}
changes only \texttt{Read}: instead of \texttt{$\Xi$ ConnReg}, it sweeps a
lapsed connection out of \texttt{registered} exactly as
\texttt{HubClientRegistry.\_sweep\_locked} does today, and leaves
\texttt{owner} untouched -- \texttt{named\_sessions} discards the swept
set's identity the moment it returns, so nothing downstream can release
what it just removed.

\begin{verbatim}
  % fixed (Read, this spec):
  \Xi ConnReg

  % buggy (Read, the H1 control): the same removal Depart performs, but
  % with no release, fired from a query instead of the coordinator
  registered' = registered \ {c? | c? : registered & lease(c?) = llapsed}
  owner' = owner
\end{verbatim}

\paragraph{The trace.}
\begin{enumerate}
  \item \texttt{Connect}(c1, i1): \texttt{registered = \{c1\}}.
  \item \texttt{Claim}(c1, s1): \texttt{owner(s1) = conn(c1)}.
  \item \texttt{TransportDies}(c1): \texttt{transport(c1) = tdown}.
  \item \texttt{LeaseLapse}(c1): \texttt{lease(c1) = llapsed}.
  \item \texttt{Read} -- buggy: \texttt{registered = \{\}} (c1 swept), but
    \texttt{owner(s1) = conn(c1)} is left exactly as it was. In the fixed
    spec, \texttt{Read} is \texttt{$\Xi$ ConnReg} and this step does
    nothing at all -- \texttt{registered} still holds \texttt{c1}.
  \item \texttt{Connect}(c2, i1): a live reconnect under the same
    identity.
  \item \texttt{Claim}(c2, s1): under the buggy variant,
    \texttt{owner(s1) = conn(c1)} and \texttt{$c1 \neq c2$}, so
    \texttt{Claim} is disabled -- the reconnecting client is shadowed by
    a connection a read, not even a write, silently removed. Worse than
    round 1's shape: no operation in this variant can ever release
    \texttt{s1}, because c1 is no longer \texttt{registered} for
    \texttt{Depart} or \texttt{Claim}'s sweep to find. In the fixed spec,
    step 5 changed nothing, so \texttt{c1} is still \texttt{registered}
    with a lapsed lease; a subsequent \texttt{Depart}(c1) or \texttt{Claim}
    by any third connection releases \texttt{s1} normally, and
    \texttt{Claim}(c2, s1) succeeds once it does.
\end{enumerate}

\begin{verbatim}
  # must be FOUND against the H1 control, NOT found against the fixed spec
  probcli <spec>.tex -model_check -nodead \
    -goal "#s.(#c.(s:SCENE & c:CONN & owner(s) = conn(c) &
                   not(c : registered)))" \
    -p DEFAULT_SETSIZE 3
\end{verbatim}

\subsection{H2 --- graceful-departure-without-release}
\label{sec:fidelity-h2}

\texttt{docs/connection\_lease\_reaping\_no\_release\_buggy.tex} changes
only \texttt{Depart}: it still removes \texttt{c?} from
\texttt{registered}, but \texttt{owner} is left untouched -- the literal
current behaviour of \texttt{HubDisplay.drop\_connection}, whose own
docstring says the departing connection's roots ``stay owned by its id.''
Round 1 kept this filename for a variant of the \emph{lease-lapse} path;
this round retargets the same filename to the operation the review
actually found broken, \texttt{Depart}, and the trace below never lapses
a lease at all -- \texttt{c1}'s transport never dies, matching the review
finding that this is the ordinary, live-session departure path, not a
timeout path.

\begin{verbatim}
  % fixed (Depart, this spec):
  registered' = registered \ {c?}
  owner' = owner (+) ({s : SCENE | owner(s) = conn(c?)} * {unowned})

  % buggy (Depart, the H2 control): the literal shipped drop_connection
  registered' = registered \ {c?}
  owner' = owner
\end{verbatim}

\paragraph{The trace.}
\begin{enumerate}
  \item \texttt{Connect}(c1, i1): \texttt{registered = \{c1\}},
    \texttt{transport(c1) = tup} throughout -- this connection's
    transport never dies.
  \item \texttt{Claim}(c1, s1): \texttt{owner(s1) = conn(c1)}.
  \item \texttt{Depart}(c1) -- buggy, an ordinary graceful disconnect,
    no lease condition required: \texttt{registered = \{\}}, but
    \texttt{owner(s1) = conn(c1)} unchanged.
  \item \texttt{Connect}(c2, i1); \texttt{Claim}(c2, s1) is disabled --
    the reconnecting client is shadowed by a connection that left
    cleanly, with a live transport, not a crashed one.
\end{enumerate}

\begin{verbatim}
  # must be FOUND against the H2 control, NOT found against the fixed spec
  probcli <spec>.tex -model_check -nodead \
    -goal "#s.(#c.(s:SCENE & c:CONN & owner(s) = conn(c) &
                   not(c : registered)))" \
    -p DEFAULT_SETSIZE 3
\end{verbatim}

\subsection{H3 --- self-reap-without-renew}
\label{sec:fidelity-h3}

\texttt{docs/connection\_lease\_reaping\_self\_reap\_buggy.tex} changes
only \texttt{Claim}: it omits the self-renewal step and the
\(c2 \neq c?\) exclusion from the sweep it performs, so the caller is
swept by its own, ordinary sweep exactly as
\texttt{HubDisplay.\_reap\_and\_release} sweeps every session including
its own caller's, because \texttt{apply} never calls anything that
renews the caller first.

\begin{verbatim}
  % fixed (Claim, this spec): renew self first, sweep only OTHERS
  lease' = lease (+) {c? |-> llive}
  registered' = registered \ {c2 : CONN | c2 : registered & c2 /= c? &
                                            lease(c2) = llapsed}

  % buggy (Claim, the H3 control): no self-renewal, no self-exclusion --
  % the caller is swept exactly like anyone else
  lease' = lease
  registered' = registered \ {c2 : CONN | c2 : registered &
                                            lease(c2) = llapsed}
\end{verbatim}

\paragraph{The trace.} \texttt{transport(c1) = tup} at every step -- this
connection never disconnects, never crashes, and never has its transport
die. It is reaped purely because its own writes never renew it.

\begin{enumerate}
  \item \texttt{Connect}(c1, i1): \texttt{registered = \{c1\}},
    \texttt{lease(c1) = llive}.
  \item \texttt{Claim}(c1, s1) -- the initial show: guard passes
    (\texttt{owner(s1) = unowned}); \texttt{owner(s1) = conn(c1)}. Under
    the buggy variant, \texttt{lease(c1)} is left exactly as it was --
    still \texttt{llive} at this point, so no divergence from the fixed
    spec is visible yet.
  \item \texttt{LeaseLapse}(c1): enabled because \texttt{lease(c1) =
    llive} and (Section~\ref{sec:leaselapse}) this operation no longer
    requires a dead transport -- \texttt{lease(c1) = llapsed}. Nothing
    renewed \texttt{c1} between steps 2 and 3; this is the patching
    connection ageing past its own TTL with only writes, no explicit
    \texttt{identify}, in between.
  \item \texttt{Claim}(c1, s1) again -- a later patch, buggy: the sweep
    set is \(\{c2 \mathbin{|} c2 \in registered \land lease~c2 = llapsed\}
    = \{c1\}\), since nothing excludes the caller. \texttt{$registered' =
    registered \setminus \{c1\} = \emptyset$}: \emph{c1 removes itself.}
    \texttt{owner} is first released for every scene c1 owned
    (\texttt{$owner(s1) \mapsto unowned$}), then the claim in this same
    call re-grants \texttt{s1} to \texttt{conn(c1)} -- so \texttt{owner(s1)
    = conn(c1)} again, but \texttt{$c1 \notin registered$}: I5/I6's
    negation is satisfied by this single step. Under the fixed spec, the
    same step's sweep set excludes \texttt{c1} by construction, so
    \texttt{registered' = \{c1\}} -- c1 stays registered, and the renewal
    in step 1 of \texttt{Claim} means this exact sequence could not have
    reached \texttt{lease(c1) = llapsed} again without a fresh, unrenewed
    gap.
\end{enumerate}

ProB's search confirms an even shorter witness than the hand-derived one
above: \texttt{Connect}(c1, i1); \texttt{LeaseLapse}(c1); \texttt{Claim}(c1,
s1) -- three steps, no successful claim needed before the lapse at all,
because \texttt{LeaseLapse} only needs \texttt{lease(c1) = llive}, which
\texttt{Connect} already establishes. The four-step version above is kept
in prose because it names the scenario the review actually reported -- an
initial show followed by a later patch -- but the mechanism is identical:
whichever step first exposes a lapsed lease at the caller's own next
\texttt{Claim}, the buggy sweep catches it.

\begin{verbatim}
  # must be FOUND against the H3 control, NOT found against the fixed spec
  probcli <spec>.tex -model_check -nodead \
    -goal "#s.(#c.(s:SCENE & c:CONN & owner(s) = conn(c) &
                   not(c : registered) & transport(c) = tup))" \
    -p DEFAULT_SETSIZE 3
\end{verbatim}

\subsection{Cascade-incomplete --- round 3's own hole, reproduced}
\label{sec:fidelity-cascade-incomplete}

\texttt{docs/connection\_lease\_reaping\_cascade\_incomplete\_buggy.tex}
changes only \texttt{TimedReapEnd} and \texttt{Claim}: both omit the three
new \(\oplus\) terms, leaving \texttt{hasSubs}/\texttt{hasWriter}/
\texttt{inboxNonEmpty} exactly as they were rather than clearing them to
\texttt{fclear}, while \texttt{Depart} is left correct -- the literal
shape of the gap Section~1 of the design document reports:
\texttt{HubDisplay.drop\_connection} already reaches all four cascade legs
once fixed, but \texttt{\_depart\_lapsed} (\texttt{TimedReap}'s
realization) and \texttt{apply}'s embedded sweep (\texttt{Claim}) never
called \texttt{Hub.on\_disconnect} or the transport sink at all.

\begin{verbatim}
  % fixed (TimedReapEnd, this spec):
  hasSubs' = hasSubs (+) {c? |-> fclear}
  hasWriter' = hasWriter (+) {c? |-> fclear}
  inboxNonEmpty' = inboxNonEmpty (+) {c? |-> fclear}

  % buggy (TimedReapEnd, the control): the lock releases, but the
  % cascade tail this round adds never runs
  hasSubs' = hasSubs
  hasWriter' = hasWriter
  inboxNonEmpty' = inboxNonEmpty
\end{verbatim}

\paragraph{The trace}, matching the design document's own sketch (Section~5)
almost verbatim.
\begin{enumerate}
  \item \texttt{Connect}(c1, i1): \texttt{registered = \{c1\}},
    \texttt{hasWriter(c1) = fset} (\texttt{Connect}'s own effect).
  \item \texttt{Activate}(c1): the abstract subscribe/writer-bind/inbox-write
    step -- \texttt{hasSubs(c1) = hasWriter(c1) = inboxNonEmpty(c1) = fset}.
  \item \texttt{TransportDies}(c1): \texttt{transport(c1) = tdown}.
  \item \texttt{LeaseLapse}(c1): \texttt{lease(c1) = llapsed}.
  \item \texttt{TimedReapBegin}(c1): \texttt{registered = \{\}} (c1
    removed), \texttt{owner} released, \texttt{cascadeFor = \{c1\}}. The
    three flags are untouched by \texttt{TimedReapBegin} in \emph{either}
    spec -- this is not yet where the two diverge.
  \item \texttt{TimedReapEnd}(c1) -- buggy: \texttt{cascadeFor = \{\}}
    again, but \texttt{hasSubs(c1) = hasWriter(c1) = inboxNonEmpty(c1) =
    fset} remain exactly as step 2 left them. \texttt{c1} is outside
    \texttt{registered} and outside \texttt{cascadeFor}, yet all three
    flags still read \texttt{fset}: I7's negation is satisfied. Under the
    fixed spec, the same step clears all three to \texttt{fclear}, so I7
    holds at every subsequent state too.
\end{enumerate}

\begin{verbatim}
  # must be FOUND against this control, NOT found against the fixed spec
  probcli <spec>.tex -model_check -nodead \
    -goal "#c.(c:CONN & not(c : registered) & not(c : cascadeFor) &
               (hasSubs(c) = fset or hasWriter(c) = fset or
                inboxNonEmpty(c) = fset))" \
    -p DEFAULT_SETSIZE 3
\end{verbatim}

Confirmed: \texttt{FOUND} against this control at \texttt{DEFAULT\_SETSIZE
3} (the 8-step trace above, found quickly); \texttt{NOT found} against the
fixed spec at \texttt{DEFAULT\_SETSIZE 2} (4{,}705 states, exhaustive --
this is I7's own check, Section~\ref{sec:verify}).

The identical control also demonstrates \texttt{Claim}'s half of the same
gap: replace steps 3--6 above with a second connection \texttt{c2} calling
\texttt{Claim}(c2, s2) once \texttt{c1}'s lease has lapsed (no
\texttt{TransportDies} needed -- \texttt{LeaseLapse} no longer requires a
dead transport). The buggy \texttt{Claim} sweeps \texttt{c1} out of
\texttt{registered} and releases its ownership exactly as the fixed spec
does, but -- with the same three \(\oplus\) terms omitted from its sweep --
leaves \texttt{hasSubs(c1)}/\texttt{hasWriter(c1)}/\texttt{inboxNonEmpty(c1)}
at whatever \texttt{Activate} set them to. The same goal is \texttt{FOUND}
by this route too, with no \texttt{TimedReapBegin}/\texttt{TimedReapEnd}
in the trace at all -- isolating \texttt{Claim}'s share of the hole from
\texttt{TimedReap}'s.

\subsection{Reconnect-races-reap --- the register\_client-vs-reap hazard
(design document Section~4)}
\label{sec:fidelity-reconnect-race}

\texttt{docs/connection\_lease\_reaping\_reconnect\_race\_buggy.tex}
changes only \texttt{Connect} and \texttt{Renew}: both drop the
\texttt{$cascadeFor = \emptyset$} guard this round adds, while every other
schema -- including the \texttt{TimedReapBegin}/\texttt{TimedReapEnd}
split itself -- is identical to the fixed spec. This is the real-code
counterpart of \texttt{register\_client}/\texttt{identify\_client} not
taking \texttt{StoreLock}: the split already exists in the fixed real
system (the multi-statement critical section \texttt{\_depart\_lapsed} and
its cascade tail run through), so what distinguishes fixed from buggy here
is purely whether the one caller that currently ignores the lock keeps
ignoring it.

\begin{verbatim}
  % fixed (Connect, this spec):
  cascadeFor = {}
  registered' = registered \/ {c?}
  ... hasWriter' = hasWriter (+) {c? |-> fset} ...

  % buggy (Connect, the control): the StoreLock fix omitted -- no guard
  % against a departure coordinator's cascade being open
  registered' = registered \/ {c?}
  ... hasWriter' = hasWriter (+) {c? |-> fset} ...
\end{verbatim}

\paragraph{The trace}, the design document's own Section~5 sketch realized
step for step.
\begin{enumerate}
  \item \texttt{Connect}(c1, i1): \texttt{registered = \{c1\}},
    \texttt{hasWriter(c1) = fset}, \texttt{lease(c1) = llive}.
  \item \texttt{LeaseLapse}(c1): \texttt{lease(c1) = llapsed} -- an
    ordinary lease timeout; no crash, no disconnect.
  \item \texttt{TimedReapBegin}(c1): \texttt{registered = \{\}} (c1
    removed), \texttt{cascadeFor = \{c1\}}. \texttt{hasWriter(c1)} is
    still \texttt{fset} -- untouched by \texttt{TimedReapBegin} in either
    spec, exactly as Section~\ref{sec:fidelity-cascade-incomplete}'s
    control shows, but here that is not yet the bug: this connection's
    cascade has genuinely not finished.
  \item \texttt{Connect}(c1, i1) again -- \textbf{buggy}: a same-identity
    reconnect, deriving the identical \texttt{c1} (design document
    Section~1.1). The fixed spec's guard, \texttt{$cascadeFor =
    \emptyset$}, would refuse this step outright, since
    \texttt{cascadeFor = \{c1\}}; the buggy control has no such guard, so
    it fires: \texttt{registered = \{c1\}} again, \texttt{hasWriter(c1) =
    fset} again (\texttt{Connect}'s own effect, representing
    \texttt{ensure\_writer} installing what it believes is a fresh
    binding), \texttt{cascadeFor} untouched by \texttt{Connect} in either
    spec -- still \texttt{\{c1\}}, since only \texttt{TimedReapEnd}
    clears it.
  \item \texttt{TimedReapEnd}(c1): enabled because \texttt{cascadeFor =
    \{c1\}} still holds -- the stale reap's cascade tail, decided back at
    step 3, finally lands. \texttt{hasWriter(c1) = fclear},
    \texttt{cascadeFor = \{\}}. The reconnect's fresh writer binding from
    step 4 is retroactively destroyed by a reap that was already stale by
    the time its cascade tail ran -- exactly the design document's own
    description, ``the two operations' effects apply out of the order
    their real-world causes occurred in.''
\end{enumerate}

After step 5: \texttt{c1 \(\in\) registered} (from the buggy step-4
reconnect) \emph{and} \texttt{hasWriter(c1) = fclear} (from step 5) -- a
live, just-reconnected session with no writer bound, silently. Under the
fixed spec, step 4 is refused; \texttt{Connect}(c1, i1) can only succeed
\emph{after} \texttt{TimedReapEnd}(c1) has cleared \texttt{cascadeFor},
and by then \texttt{Connect}'s own unconditional
\texttt{$hasWriter' = hasWriter \oplus \{c? \mapsto fset\}$} is the last
word: \texttt{hasWriter(c1) = fset} once the reconnect completes, with
nothing left in \texttt{cascadeFor} to overwrite it afterward.

\begin{verbatim}
  # must be FOUND against this control, NOT found against the fixed spec
  probcli <spec>.tex -model_check -nodead \
    -goal "#c.(c:CONN & c : registered & hasWriter(c) = fclear)" \
    -p DEFAULT_SETSIZE 3
\end{verbatim}

Confirmed: \texttt{FOUND} against this control at
\texttt{DEFAULT\_SETSIZE 3} (a 10-step trace, found quickly -- locating one
witness needs no exhaustive search); \texttt{NOT found} against the fixed
spec at \texttt{DEFAULT\_SETSIZE 2} (4{,}705 states, exhaustive).
\texttt{DEFAULT\_SETSIZE 3} against the fixed spec was attempted and did
not complete this round (Section~\ref{sec:verify}'s Deadlock paragraph
records the same state-space growth); \texttt{DEFAULT\_SETSIZE 2}'s
exhaustive result is what the \texttt{NOT found} claim above rests on.

\end{document}
